\section{The large-prime field on the selected family}

% BEGIN J80H1
\paragraph{J80H1. The rational local inputs on the specified period rings.}
Use the Bloch--Kato, Tate duality and Euler characteristic fields of
$\mathrm{Hypotheses}_2$ at $K$. For the first field use the tie of
$D$ at this prime. Applied to the underlying $\mathbb Q_p$-space,
its dimension formula computes the finite subspace by adding
$\dim H^0$ to $\dim D_{\rm dR}/\mathrm{Fil}^0$. A crystalline
representation here is de Rham by W79H1. Restricting scalars
introduces the same factor $[L:\mathbb Q_p]$ in all three
dimensions: the period tensors are over $\mathbb Q_p$ in both
conventions, and invariants, coboundaries, the comparison defining
finiteness, and the filtration are stable under $L$. In the split
algebra $L\otimes K$ separate the embedding idempotents. The dimension
of each quotient by $\mathrm{Fil}^0$ is the sum of its graded
dimensions in negative degrees, by exhaustive finite filtrations.
This is the jump formula used in T80B7. The integral finite condition
is saturated because it is the inverse image of an $L$-subspace in
rational cohomology. For the trivial line its finite dimension is
one. The unramified degree cocycle is finite: its triangular
representation on two basis vectors has a stable lattice, is
unramified, and the Witt matrix lifting argument of S80J4 provides
a crystalline splitting after extending scalars to the period ring.
The degree is nonzero on Frobenius. Continuous unramified
homomorphisms to $L$ are precisely its $L$-multiples by density
of integral Frobenius powers. Thus the full rational unramified
clause also holds. Tate duality and the Euler formula apply directly
to the finite underlying $p$-groups with their continuous actions.

For clarity, in the crystalline extension of T80B8 the filtration
on the quotient is strict. Take invariant comparison bases on the
three terms. The sequence of crystalline modules is exact, since
after scalar extension by the crystalline comparison it is the
original exact sequence tensored with the period ring; on each
letter extension of scalars from a field to a nonzero ring reflects
exactness. The same holds for de Rham modules. Filtration on the
kernel is intersection since the filtration on the representation
tensor is induced by $B_{\rm dR}^+$. The filtration induced on the
quotient is contained in its own period filtration. To see equality,
at each embedding the valuation of the period determinant is the
sum of the jumps: W79H2--W79H3 give the integral period basis
$t^{-u_j}b_j$ for a filtration-adapted crystalline basis $b_j$, up
to a unit determinant. Determinants in the exact sequence multiply,
using compatible bases on the representation and invariant sides.
Algebraic changes of basis are units of the period valuation ring.
Thus the sum of the jumps in the middle term equals the two outer
sums. In a filtered exact sequence with the induced quotient
filtration the graded dimensions add by construction. A proper
inclusion in any step of the two quotient filtrations would give a
strict inequality of their jump sums: subtract successive dimensions
and telescope on their common finite filtration range. They are
therefore equal at every step and label. This justifies the use of
strictness on precisely the extensions needed here.
% END J80H1

% BEGIN J80H2
\paragraph{J80H2. From all three local conclusions to the large-prime estimate.}
Fix a family as in D79G8. Take the maximum of the cutoffs for the
jump readings, position tables and parameter inequalities in S80J1--S80J5.
These depend on the family before the choice of any localization.
At any later member satisfying its native facts, take an arbitrary
localization, good family $R$, point $\pi$, complement $\rho$ and
place $w\mid\ell$ under the large-prime guards of D79G8. S80J5
produces the good unramified local layer with its coefficient
enlargement. Use the compatible local records from L78W13 on this
layer, available from the Bhatt--Scholze hypothesis. R80D6 provides
the residual presentation for each residue field. W80N5 and M80R4
turn it into the full residual realization. Together with J80H1 and
the residual and finite-coefficient arguments of T80B1--T80B10, this
gives the three native cohomological tests on these very local data.
The shifted jumps and characteristic-polynomial congruences for all
points required by the tests are supplied by S80J6. Apply C80L6
and then L79C1--L79C2. At $w$ each original rank-block cut generator
is Bloch--Kato finite, each character-ratio generator with index
$2\le c\le t$ is a coboundary, and every first-character cut is
unramified. All choices have been made inside the quantifiers of
D79G8. The congruence bound at the member follows from that theorem
with the global Artin reciprocity input in $\mathrm{Hypotheses}_2$.
The rank-block local inequality uses zero additional error; the
finite-part term in the inequality can be nonzero.

Finally take the common constant of Y79T4 and increase its cutoff
to include the one just obtained. For every member satisfying its
scale and finiteness guards keep the very hub of Y79T4, with its
prescribed exceptional set. G79E10 applies the congruence estimate
there and gives $\mathrm{Concl.large}$. On this one hub we have
\[
 \begin{gathered}
 \mathrm{pos},\ \mathrm{pass},\ \mathrm{mult},\ \mathrm{rank},
 \mathrm{kill},\ \mathrm{num},\ \mathrm{bb},\\
 \mathrm{exc},\ \mathrm{typ},\ \mathrm{par},\ \mathrm{small},
 \mathrm{large}.
 \end{gathered}
\]
The original denominator family conclusion still requires the
middle-range, small-prime correction and ramified-prime fields.
% END J80H2
